%!TEX root = ../hutbthesis_main.tex

\chapter{. 系统设计与实现}\label{ux7cfbux7edfux8bbeux8ba1ux4e0eux5b9eux73b0}

\section{实验环境与整体框架}\label{ux5b9eux9a8cux73afux5883ux4e0eux6574ux4f53ux6846ux67b6}

（1）软硬件运行环境

实验用的是AMD Radeon（TM）图形处理器，搭配NVIDIA GeForce GTX
1650独立显卡，内存为16GB，硬盘为512GB固态硬盘。该配置可以满足肌肉骨骼动力学仿真的需要，可以满足深度强化学习训练的需求，还可以处理大量的运动数据。

软件环境层面，系统运行于 Windows 11 操作系统，基于 Python 3.9
语言开发，使用 PyCharm
作为集成开发工具。仿真部分利用Kinesis平台创建肌肉骨骼模型，使用MuJoCo
\subsection{物理引擎做高精度动力学仿真，深度学习模块用PyTorch}
\subsection{框架开发，整体环境稳定，可以支持复杂的运动控制算法的开发和验证。开发环境如表2所示。}

\begin{center}\textbf{表2 开发环境}\end{center}

\begin{longtable}[]{@{}
  >{\raggedright\arraybackslash}p{(\columnwidth - 4\tabcolsep) * \real{0.1594}}
  >{\raggedright\arraybackslash}p{(\columnwidth - 4\tabcolsep) * \real{0.3463}}
  >{\raggedright\arraybackslash}p{(\columnwidth - 4\tabcolsep) * \real{0.4943}}@{}}
\toprule\noalign{}
\endhead
\bottomrule\noalign{}
\endlastfoot
开发平台 & 软件环境 & 硬件环境 \\
\multirow{3}{=}{PC} & 操作系统：Windows 11 & \multirow{3}{=}{CPU：AMD
Radeon(TM)Graphics

GPU：NVIDIA GeForce GTX 1650} \\
& 开发语言：Python \\
& IDE：pycharm \\
\end{longtable}

（2）实验参数设置

为了使实验结果可以重复、比较，所有的训练以及测试过程都使用相同的参数设置。实验随机种子为42，可以保证训练随机性可控；仿真步长取0.01s，既能保证仿真精度又能提高计算效率；批量大小取64，梯度更新稳定，训练速度快。

关键超参数中学习率设为3×10\^{}-4，适合深度网络参数更新速度；PPO裁剪系数设为0.2，限制策略更新幅度防止训练震荡；折扣因子设为0.99，平衡短期奖励和长期收益；优势估计系数设为0.95，减小价值估计偏差；最大训练步数设为200万步，保证算法充分收敛得到最优控制策略。

（3）系统整体实现框架

本系统用深度学习和物理仿真相结合的方式进行搭建，采用模块化开发的思想，把数据处理、运动学计算、策略控制、物理仿真、量化评估这五个功能统一起来，形成一个可以直接运行、可以复现、可以评价的肌肉骨骼运动模仿完整系统。系统从实现上是高度标准化的，所有的流程固定、接口一致、参数可以配置，可以稳定地产生可比的实验结果。

系统整体实现逻辑就是加载、标准化运动数据，给策略提供一个统一的参考轨迹，再用前向运动学模块将关节姿态转换成三维空间状态，Actor根据状态信息输出连续的控制信号，仿真引擎执行控制指令更新肌肉骨骼模型状态，最后由评估模块计算精度、稳定性、鲁棒性指标形成闭环。

\section{核心功能模块实现}\label{ux6838ux5fc3ux529fux80fdux6a21ux5757ux5b9eux73b0}

本研究核心功能模块分为运动数据处理模块、前向运动学模块、运动控制模块、仿真与评估模块。系统核心功能如图4-1所示。

\begin{center}
\includegraphics[width=0.92\textwidth]{images/image8.png}
\end{center}

\begin{center}\textbf{图4-1 系统核心功能示意图}\end{center}

\subsection{运动数据处理模块实现}\label{ux8fd0ux52a8ux6570ux636eux5904ux7406ux6a21ux5757ux5b9eux73b0}

运动数据处理模块把原始的运动捕获数据转变成规范的、平滑的、可以用来做策略学习的标准序列，是保证模仿准确性和实验稳定的根基。

模块统一读取运动序列文件，解析关节旋转、根节点位移、速度、角速度等信息，把不同的格式输入统一成系统可以识别的结构。

在预处理阶段先进行噪声滤波、根节点高度校正、朝向归一化、帧率统一、异常姿态剔除等处理使得所有的动作空间、时间、姿态趋于一致，为后面策略学习与评价打下良好的基础。

数据增强模块在训练的时候会随机对运动的朝向进行旋转，加上很微小的噪声，时间移动以及速度调整这些操作，从而使得有成效的样本数量变多，进而加强策略对于各种环境、各种起始状况的适应程度以及抵抗性的能力。

改进优化内容：

多线程并行加载使用多进程队列的方式并行读取、预处理运动数据，大大提高大体积数据的加载速度。

标准化预处理优化就是用滤波、对齐、归一化、异常剔除等一系列的处理方法来保证数据在同一个标准下进行学习，提高实验的可重复性。

数据增强技术通过将训练集的数据做旋转、小幅度改变、速度增大等操作来扩大样本数量，从而降低过拟合概率，提升策略的泛化程度。

\subsection{前向运动学与运动控制模块实现}\label{ux524dux5411ux8fd0ux52a8ux5b66ux4e0eux8fd0ux52a8ux63a7ux5236ux6a21ux5757ux5b9eux73b0}

前向运动学模块完成状态解算的主要工作，把关节的旋转信息转化为三维关节的位置、速度和角速度，给控制器提供可观测的状态空间。

模块采用SMPL参数化人体模型，可以进行批量、高效、多帧并行的前向运动学计算。模块运行时会自动对关节进行顺序映射，将坐标系进行转换，再做姿态插值和速度平滑，使状态信息正确、低噪、连续。输出全局位置、局部姿态、根节点状态等多级信息给Actor，给Actor赋予丰富的稳定观测输入。

改进优化内容：

批量并行计算优化，用向量化运算做多帧同步计算，提高运动学求解速度，满足大规模数据处理的要求。

关节顺序映射优化，创建SMPL和MuJoCo的关节索引映射表，自动把跨模型的关节顺序进行重排，从而达到数据无缝对接的目的。

坐标系转换用四元数和旋转矩阵来完成，不会出现欧拉角万向锁的问题，姿态计算连续稳定。

运动控制模块为系统的决策中枢，用深度强化学习把状态和动作联系起来。Actor把关节角度、速度、参考轨迹当作输入，输出平滑、有界且符合生理约束的控制信号。Critic和Actor有相似之处，利用两者的共同性即共享特征提取层得到的当前运动长期收益状态价值分数来评价运动的好坏，以此降低策略更新时出现的剧烈波动，避免训练出现过于剧烈或者缓慢的情况。

\subsection{仿真与评估模块实现}\label{ux4effux771fux4e0eux8bc4ux4f30ux6a21ux5757ux5b9eux73b0}

（一）物理仿真模块使用MuJoCo引擎创建出一个高保真的肌肉骨骼动力学环境，模拟骨骼刚体、关节约束、肌肉驱动、地面接触、摩擦和碰撞等各种物理特性，严格按照生物力学的规律。仿真时及时更新姿态、位置、速度、接触状态，对超出生理极限的动作加以限制或者停止，保证运动过程的安全、稳定、可信。

改进优化内容：

高精度物理仿真优化就是对肌肉执行器、关节限位、力矩约束进行完善的处理，使仿真更加符合真实的人体力学特性。

无头模式加速，关闭图形渲染，只保留物理计算，大大提高了训练和评估的速度，可以进行大量的重复实验。

（二）
评估与反馈模块实现，评估模块完成标准化、可重复的量化评价，在每一个运动执行过程里即时存档姿态、位置、误差以及运行帧数，最后得出诸多重要指标，给算法改良赋予客观支撑。

改进优化内容：

MPJPE：平均每关节位置误差，衡量模仿精度。
\begin{equation}
MPJPE=\frac{1}{T}\sum_{t=1}^{T}\frac{1}{J}\sum_{j=1}^{J}\left\|x_j(t)-\widehat{x}_j(t)\right\|_2
\tag{4-1}
\end{equation}

帧覆盖率：反映运动连续执行稳定性。

任务成功率：反映策略鲁棒性与工程可用性。

为确保评估结果可靠可验证，需将评估流程标准化，包括采用固定测试集、设定随机种子、统一运行参数及评估帧数，以此保障多次运行结果可重复、可对比。

\section{生物合理性约束实现}\label{ux751fux7269ux5408ux7406ux6027ux7ea6ux675fux5b9eux73b0}

为了使运动符合人体生理规律，在系统实现上把生物合理性约束嵌入到奖励函数和仿真环境中，形成了一个完整的硬约束加软惩罚的机制。

关节角度约束结合仿真限位和惩罚函数，直接限制关节活动在生理范围内，不让它随意超出。肌肉驱动的约束让肌肉的激活幅度和速度都收紧一些，这样肌肉的表现更像真实的力学反应。能量消耗约束专门针对高耗能、动作冗余或者突然变化的动作进行惩罚，推动个体选择更自然、更节能的运动方式。至于运动平滑性约束，它通过限制加速度和加加速度，用滤波手段让动作变得连续流畅，更贴近真实运动的样子。

依靠多层约束的共同作用，在保证轨迹精度的同时达到高度自然、高度可信的生物合理运动。目标函数采用多因素加权组合形式：\\
\begin{equation}
L=w_1L_{pose}+w_2L_{vel}+w_3L_{bio}+w_4R_{task}
\tag{4-2}
\end{equation}

各项含义如下：

\(L_{pose}\)：姿态与关节位置损失，保证空间轨迹与参考动作相似；

\(L_{vel}\)：速度损失，保证运动快慢、动态变化趋势与示范一致；

\(L_{bio}\)：生物合理性惩罚项，包括关节超限惩罚、肌肉异常激活惩罚、能量冗余惩罚等；

\(R_{task}\)：任务奖励，包括完成动作奖励、保持平衡奖励、无崩溃奖励等。

\section{系统运行效果测试与分析}\label{ux7cfbux7edfux8fd0ux884cux6548ux679cux6d4bux8bd5ux4e0eux5206ux6790}

\subsection{肢体运动分析}


（1）下肢运动协同性分析

下肢交替协同是评价步态生物合理性的重要指标，也是决定生成动作是否自然、运动是否流畅的关键。实验得到脚趾、膝关节、踝关节时序运动数据，计算左右侧运动相关性和空间对称误差，结果见表3。

\begin{center}\textbf{表3 下肢协同量化指标统计结果}\end{center}

\begin{longtable}[]{@{}
  >{\raggedright\arraybackslash}p{(\columnwidth - 4\tabcolsep) * \real{0.3333}}
  >{\raggedright\arraybackslash}p{(\columnwidth - 4\tabcolsep) * \real{0.3334}}
  >{\raggedright\arraybackslash}p{(\columnwidth - 4\tabcolsep) * \real{0.3334}}@{}}
\toprule\noalign{}
\endhead
\bottomrule\noalign{}
\endlastfoot
评估指标 & 数值 & 说明 \\
左右脚交替相关系数 & 0.764 & 交替时序一致性高 \\
膝角度相关系数 & 0.139 & 屈伸相位差合理 \\
踝高度对称误差 /m & 0.0303 & 空间对称性良好 \\
\end{longtable}

从图4-2的轨迹曲线可以直观看出，左右脚趾的前后位移呈现出非常明显的反向交替运动。这种特征和0.764的高相关系数相吻合，说明步态中左右脚交替节奏稳定，没有明显时序偏移。再看左右膝关节的角度曲线，周期性的屈伸变化非常清楚，两条曲线的相位差和真实人体步态的膝关节运动模式几乎一致。这也解释了为什么膝角度的相关系数偏低：不是因为协同性差，而是因为这种相位差本身就是符合人体生物力学的自然设计。至于左右脚踝的高度曲线，每次峰值都对应该步态周期中的抬脚时刻，而且两条曲线的峰值误差极小，和0.0303m的对称误差数据相印证，进一步证明了动作空间高度的对称性。证明了动作空间高度的对称性，整体下肢运动完全符合人体步态的生物力学特征。

\begin{center}
\includegraphics[width=0.92\textwidth]{images/image9.png}
\end{center}

\begin{center}\textbf{图4-2 下肢关节运动时序轨迹可视化曲线}\end{center}

（2）上半身运动稳定性分析

人体行走时上半身应有小幅、对称、平稳的摆动，从而保证躯干的稳定，减少运动耗能，这是评价动作生物合理性的又一重要方面。实验得到腕关节、肘关节、肩关节时序运动数据，对生成动作上半身运动特征进行可视化分析，结果如图
4-3 所示。

\begin{center}
\includegraphics[width=0.92\textwidth]{images/image10.png}
\end{center}

\begin{center}\textbf{图4-3 上半身关节运动时序轨迹可视化曲线}\end{center}

从图4-3的轨迹曲线来看，左右腕关节在前后方向上有明显的对称且幅度较小的摆动，左右摆臂的相关系数是-0.683，这正好符合人在行走时两臂反向摆动的规律。再看左右肘关节的角度曲线，几乎完全重合，相关系数高达0.994，说明肘炎关节的运动在两侧极为一致。左右肩的高度对称误差只有0.0241米，进一步证明上半身的运动在空间上非常对称。整体来看，上半身的运动幅度不大，波动也很平稳，并且左右对称。这些结果直接反映出算法对躯干的稳定控制能力很强，生成的运动自然、节能，完全符合人体自身的平衡调节机制。

\subsection{奖励函数优化效果分析}

为了检验新增生物约束奖励的优化效果，我设计了两个对照实验。基础奖励组只考虑了轨迹跟踪、速度约束、直立姿态和运动能耗这四项基础奖励。而优化奖励组在此基础上，又加入了关节角度约束、肌肉驱动代理和运动平滑性这三项与生物合理性相关的奖励。图4-4展示了优化奖励中新加入的生物合理性代价的分布情况。至于不同动作类型在奖励优化前后的表现差异，可以直接看图4-5。

\begin{center}
\includegraphics[width=0.92\textwidth]{images/image11.png}
\end{center}

\begin{center}\textbf{图4-4 奖励组成分析图}\end{center}

\begin{center}
\includegraphics[width=0.92\textwidth]{images/image12.png}
\end{center}

\begin{center}\textbf{图4-5 不同动作类型奖励函数优化前后性能对比}\end{center}

实验结果表明，新增的三个生物约束都提升了运动性能。关节角度约束的归一化代价最低，直接说明它在限制关节运动范围方面最有效。肌肉驱动代理、能耗和平滑性约束的代价也都合理，优化奖励不仅保证了基本跟踪性能，还让动作更符合生物逻辑。

看动作分组的数据，后退、左转、右转、环形行走、前向行走这五类典型动作，优化后平均奖励分数都比优化前高。加上生物约束奖励后，算法在各种动作上表现都不错，泛化能力也得到了验证。

\subsection{算法量化对比实验}


（一）实验设置

本次实验使用标准测试集和统一评价协议来保证实验结果的公平性、可重复性、可比性。

选择出50条标准的下肢运动序列作为测试数据集，包含直线行走、原地转身、侧向跨步、环形行走、变速走、小步调整等日常活动的动作类型比较平衡、难度逐渐变化，可以很好地反映该算法对于常见和复杂动作的处理能力。代码运行测试如图4-6所示。

\begin{center}
\includegraphics[width=0.92\textwidth]{images/image13.png}
\end{center}

\begin{center}\textbf{图4-6 代码运行测试图}\end{center}

对比算法：

基准算法是标准近端策略优化（PPO）加上对角高斯策略（假设动作维度独立）以及基础奖励函数（只包含位置跟踪、速度跟踪、姿态直立、能量消耗这四个项目）。

本文改进算法为PPO基础框架、Lattice全协方差策略网络、混合专家(MoE)结构、生物增强奖励函数（增加关节角度约束、肌肉驱动约束、运动平滑性约束三项）。

评价指标选择MPJPE、帧覆盖率、任务成功率三项行业公认的指标，分别对应模仿精度、运动稳定性、任务鲁棒性这三个主要方面，从而达到对算法性能的全面、客观、量化评价的目的。

（二）评估指标与实验方案

本文使用组合式的量化评价体系，选择三项行业通用的核心指标，从精度、稳定性和鲁棒性三个方面综合评价算法的性能。

平均每关节位置误差（MPJPE）用来衡量预测关节和真实关节三维位置的平均欧式距离，单位为cm，数值越小说明模仿精度越高。

帧覆盖率，算法在不崩溃、不发生异常终止的情况下完成的帧数占总帧数的比例，可以体现运动的连续稳定程度。

任务成功率为多次测试中完成所有动作并且没有跌倒、关节超限、模型穿模等异常情况出现的序列所占比例，体现的是策略的鲁棒性和工程可用性。

本文实验严格遵守标准化、可重复的原则，具体的方案如下所示：

测试数据集由已经标准化过的100个下肢动作序列组成，包含行走、跨步、转身、慢走等不同的动作

重复实验次数为每组实验独立重复 5
次，取指标的平均值作为最终结果，减小随机误差的影响；

评估流程统一使用 Kinesis 平台内置的 imitation eval
评估流程，自动计算并记录
MPJPE、帧覆盖率、任务成功率，保证评估过程客观、透明、可追溯。

采用固定实验参数、统一测试数据集、标准化评价流程的方法来消除实验误差，保证结果可以重复、可以比较、可以验证，给肌肉骨骼运动控制算法的性能评价提供一个可靠的范式。

（三）量化对比结果

\textbf{表 4 基准算法与本文算法性能对比}

\begin{longtable}[]{@{}
  >{\raggedright\arraybackslash}p{(\columnwidth - 6\tabcolsep) * \real{0.3439}}
  >{\raggedright\arraybackslash}p{(\columnwidth - 6\tabcolsep) * \real{0.1614}}
  >{\raggedright\arraybackslash}p{(\columnwidth - 6\tabcolsep) * \real{0.2474}}
  >{\raggedright\arraybackslash}p{(\columnwidth - 6\tabcolsep) * \real{0.2474}}@{}}
\toprule\noalign{}
\endhead
\bottomrule\noalign{}
\endlastfoot
算法 & MPJPE(cm) & 帧覆盖率 (\%) & 任务成功率 (\%) \\
基准算法（标准 PPO） & 28.7 & 81.2 & 72.4 \\
本文改进算法 & 21.3 & 93.5 & 84.8 \\
提升幅度 & -25.8\% & +15.2\% & +17.1\% \\
\end{longtable}

如表4所示基准算法的平均关节位置误差为28.7cm，存在明显的轨迹偏差；本文改进算法MPJPE降低到21.3cm，误差降低了25.8\%，关节轨迹跟踪精度提高了很多，姿态贴合度也更高。

基准算法帧覆盖率是81.2\%，运动时姿态会突然改变、失衡抖动、中途崩溃等会随机发生，连续执行能力差；本文改进算法帧覆盖率为93.5\%，提高了15.2\%，运动连续性与稳定性更好，动作崩溃、抖动等得到很好地抑制。

基准算法任务成功率为72.4\%，转弯、跨步等复杂动作不能很好地适应，姿态超限、平衡失稳容易造成任务失败，本文改进后的算法任务成功率提高到84.8\%，提高了17.1\%，对各种动作的适应性明显提高，在复杂环境下任务完成的可靠性也更高。

（1）MPJPE深度分析

MPJPE，也就是平均每关节位置误差，其实就是评判动作模仿准确度的关键指标。数值越低，生成的动作和参考动作之间的空间相似度就越高，几何上的偏差也更小。基准算法的整体
MPJPE 达到 28.7
厘米，这说明生成动作和标准参考轨迹之间存在明显的系统性偏差。本文分析了简单走路、转身和变速情况下的动作精度。我们重复实验，取均值进行对比，具体结果见表5。

\textbf{表 5 不同的动作MPJPE对比}

\begin{longtable}[]{@{}
  >{\raggedright\arraybackslash}p{(\columnwidth - 6\tabcolsep) * \real{0.2978}}
  >{\raggedright\arraybackslash}p{(\columnwidth - 6\tabcolsep) * \real{0.2021}}
  >{\raggedright\arraybackslash}p{(\columnwidth - 6\tabcolsep) * \real{0.2501}}
  >{\raggedright\arraybackslash}p{(\columnwidth - 6\tabcolsep) * \real{0.2501}}@{}}
\toprule\noalign{}
\endhead
\bottomrule\noalign{}
\endlastfoot
算法 & 简单走路(cm) & 转身(cm) & 变速(cm) \\
基准算法（标准 PPO） & 24.3 & 29.8 & 35.4 \\
本文改进算法 & 17.6 & 21.5 & 25.8 \\
对比提升 & 7.7cm & 8.3cm & 9.6cm \\
\end{longtable}

结论：从不同动作的MPJPE对比结果来看，基准算法在简单走路、转身、变速这三类典型动作上的误差都比较高。动作越复杂，误差增长越快，最多能达到35.4厘米。传统算法适应动态运动的能力明显不足，动作难度一上升，轨迹偏差就变得更严重。我们改进后的算法在三大动作中都大幅降低了误差：简单走路降到17.6厘米，转身21.5厘米，变速25.8厘米,比基准算法分别少7.7厘米、8.3厘米和9.6厘米。实验结果很清楚地说明，改进算法能稳定优化不同复杂度的动作，对动态运动的适应远胜传统算法，大大提升了肌肉骨骼运动模仿的精准度和可靠性。

（2）运动稳定性（帧覆盖率）深度分析

帧覆盖率直接反映了动作连续性和抗崩溃能力，这其实是生物合理性最本质的指标。基准算法覆盖率只有81.2\%，也就是说，差不多有五分之一的动作序列在执行时会坍塌、失衡、穿模，或者因关节超限而被迫终止。。我们提出的改进算法把覆盖率提升到了93.5\%，异常终止的比例降到了6.5\%。运动的连贯性和整体稳定性明显更好。表6对比了三种常见的崩溃序列问题。

\textbf{表 6 崩溃序列性能对比}

\begin{longtable}[]{@{}
  >{\raggedright\arraybackslash}p{(\columnwidth - 6\tabcolsep) * \real{0.3139}}
  >{\raggedright\arraybackslash}p{(\columnwidth - 6\tabcolsep) * \real{0.1860}}
  >{\raggedright\arraybackslash}p{(\columnwidth - 6\tabcolsep) * \real{0.2501}}
  >{\raggedright\arraybackslash}p{(\columnwidth - 6\tabcolsep) * \real{0.2501}}@{}}
\toprule\noalign{}
\endhead
\bottomrule\noalign{}
\endlastfoot
算法 & 关节超限(\%) & 失衡跌倒(\%) & 轨迹漂移(\%) \\
基准算法（标准 PPO） & 41.2 & 34.8 & 24.0 \\
本文改进算法 & 7.8 & 4.9 & 2.3 \\
对比提升 & 33.4\% & 29.9\% & 21.7\% \\
\end{longtable}

对于不同动作的稳定性提升效果差异显著，本文对简单行走、原地转身和变速行走三个动作来进行对比分析，结果如表7所示。

\textbf{表 7 不同动作帧覆盖率对比}

\begin{longtable}[]{@{}
  >{\raggedright\arraybackslash}p{(\columnwidth - 6\tabcolsep) * \real{0.3139}}
  >{\raggedright\arraybackslash}p{(\columnwidth - 6\tabcolsep) * \real{0.1860}}
  >{\raggedright\arraybackslash}p{(\columnwidth - 6\tabcolsep) * \real{0.2501}}
  >{\raggedright\arraybackslash}p{(\columnwidth - 6\tabcolsep) * \real{0.2501}}@{}}
\toprule\noalign{}
\endhead
\bottomrule\noalign{}
\endlastfoot
算法 & 简单行走(\%) & 原地转身(\%) & 变速行走(\%) \\
基准算法（标准 PPO） & 89.5 & 76.3 & 68.1 \\
本文改进算法 & 97.8 & 92.1 & 88.3 \\
对比提升 & 8.3\% & 15.8\% & 20.2\% \\
\end{longtable}

结论：看数据，基准算法的崩溃主要有三大原因：关节超限占41.2\%，失衡跌倒有34.8\%，轨迹问题也有24\%。这些几乎就是运动失败的全部来源。但用上改进算法后，这三类异常一下降到7.8\%、4.9\%、2.3\%。核心稳定性的问题基本解决了。

不同动作场景的测试也有意思。动作越难，稳定性提升越大。简单行走稳定性提升8.3\%，原地转身提升到15.8\%，变速行走提升最多，达到20.2\%。所以，改进算法在动态、变向这类高难度运动里，稳定性优势特别明显，对复杂运动需求适应得更好。

总的来说，这套改进算法绕开了传统算法的那些大坑——关节超限、姿态失衡、轨迹漂移——运动的生物力学合理性更强，原来容易崩溃、稳定性差的短板有了本质上的改善。运动的连贯性和可靠性都大幅提升。生物约束对动态和变加速动作的稳定性优化效果非常显著。简单行走的稳定性几乎满分，说明基础动作已经非常接近理想、稳定的自然运动状态了。

（3）任务成功率深度分析

本文每组实验独立重复五次，统计成功率波动其中基准算法：72.4\%±9.2\%，波动大，结果不稳定；本文算法：84.8\%±3.5\%，波动小，结果可靠。基准算法在执行过程中，失败主要集中在姿态严重变形、平衡彻底丢失、轨迹大幅偏离三类核心问题上。相比之下，本文改进算法的失败情况仅表现为轻微姿态偏差、边缘状态下的短暂失衡，完全不存在轨迹偏离这类致命性故障。

结论：改进算法从根本上杜绝了传统算法严重失效的情况,只剩下少量的微小偏差不影响整体动作,实质性地增强了任务执行的完整性和最终完成的质量。

（4）总结

从模仿精度、运动稳定性和任务鲁棒性这三项核心指标来看，这篇论文提出的改进算法在每个维度的表现都比传统基准算法强得多。动作模仿变得更准确，动态运动更稳定，遇到复杂环境时任务执行也更可靠。无论怎么比，改进后的算法都压过了基准模型一头。

这背后，Lattice全协方差策略、MoE混合专家架构，还有生物增强奖励函数这些核心模块功不可没。它们解决了传统算法精度不够、稳定性差、生物合理性不足的问题。换句话说，这套方法不仅让人体运动模仿更精确、更可靠，也让它更贴近真实生物运动。对想要实现高精度、高稳定性和强生物合理性的人体运动模仿任务来说，这套算法提供了科学、可靠，也具备工程参考价值的实现路径。
